% Root full proof and original-source audit passed before promotion.
% AI-produced by GPT-6 Astra (OpenAI, ChatGPT Pro mode), inherited effort,
% agent /root/higher_comparison, RESEARCH_SA20, 19 November 2026.
\subsection{Elementary extensions and the whole definable concordance space}

\begin{saresult}{A.133}{Proved}{Whole definable concordance spaces are invariant under elementary extension}
Let $\mathcal M\preccurlyeq\mathcal N$ be an elementary extension of
o-minimal structures in the same language, both expanding real closed
fields. Let $B$ be a closed bounded definable manifold over $\mathcal M$.
Then extension of definitions induces a weak equivalence of simplicial
groups
\[
 P_{\mathrm{def},\mathcal M}(B)
 \longrightarrow
 P_{\mathrm{def},\mathcal N}(B_{\mathcal N}).
\]
The groups use homeomorphisms of the total product with a parameter
simplex; they preserve that parameter, fix the bottom pointwise and
preserve the free top. All defining parameters in the relevant structure
are permitted. More generally, extension gives the same weak equivalence for the
simplicial automorphism group of any closed bounded definable set
with finitely many marked definable subsets, each required to be
preserved setwise or fixed pointwise.

If $B$ is given by a finite rational PL model, the square with the PL
coefficient-extension map commutes. Hence the comparison from PL is a
weak equivalence over $\mathcal M$ if and only if it is a weak equivalence
over $\mathcal N$. Corresponding injectivity, surjectivity, or isomorphism assertions
in each homotopy degree are equivalent as well.
\end{saresult}
\begin{proof}
First the extension map is well defined. Two formulas over $\mathcal M$
that define the same graph satisfy the same universal graph-equality
sentence in $\mathcal M$, hence in $\mathcal N$. Being a graph, being
bijective, preserving the parameter, satisfying the boundary conditions,
and continuity of both total maps are first-order conditions. For
continuity we use the epsilon-delta formula with all coordinates of the
total space included. They remain true under elementary extension.
Composition, inversion and affine parameter pullback are given by the
same formulas in the two structures. Thus extension is a simplicial
homomorphism.

We give the finite-coefficient continuity argument, since elementary
transfer by itself would not connect a given family to a descended one.
Fix a closed bounded definable set $D$ over $\mathcal M$ and one formula
$\varphi(a;x,y)$ with fixed parameters in $\mathcal M$. Let $T$ be the
set of tuples $a$ for which this formula defines a homeomorphism
$h_a:D\to D$ with a prescribed finite list of first-order conditions.
The set $T$ is definable, with no quantifier-elimination hypothesis on
the o-minimal theory. Put
\[
 \Gamma=\{(a,x,y):a\in T,\ x,y\in D,\ \varphi(a;x,y)\}.
\]
Apply definable Hardt triviality to $\Gamma\to T$ and refine its base
partition into cells. The exact standard input is definable Hardt triviality
\cite[Theorem~5.22, pp.~62--63]{Coste1999}. The source works over an
arbitrary real closed field in a fixed o-minimal structure
\cite[\S1.3, pp.~9--10]{Coste1999}; the triviality theorem is stated
for arbitrary definable families, not only proper maps. All this is done in $\mathcal M$. On a nonempty
cell $C$ choose $a_0\in C(\mathcal M)$. The trivialization can be
written
\[
 \theta:C\times\Gamma_{a_0}\xrightarrow{\sim}\Gamma|_C
\]
over $C$. The fibre $\Gamma_{a_0}$ is closed and bounded because it is
the graph of a continuous map on $D$.

Compose $\theta$ with the projection to $(a,x)$, obtaining a continuous
definable bijection
\[
 J:C\times\Gamma_{a_0}\longrightarrow C\times D.
\]
It is a homeomorphism. To verify inverse continuity over a possibly
non-Archimedean field, suppose it fails at $(a,x)$. Definable curve
selection supplies a definable curve in $C\times D$ tending to $(a,x)$
whose inverse stays outside some fixed neighbourhood of
$J^{-1}(a,x)$. The first coordinate of its inverse tends to $a$.
The second coordinate lies in the closed bounded set $\Gamma_{a_0}$;
definable compactness, or coordinatewise monotonicity, gives a limit
$\gamma\in\Gamma_{a_0}$. Continuity of $J$ gives
$J(a,\gamma)=(a,x)$. Bijectivity identifies $(a,\gamma)$ with
$J^{-1}(a,x)$, a contradiction. Thus the output projection of
$\theta J^{-1}$ gives joint continuous evaluation
$(a,x)\mapsto h_a(x)$ on $C\times D$.
Interchanging input and output gives joint inverse evaluation.
These maps and their continuity assertions are defined over
$\mathcal M$, so they extend to $\mathcal N$ with the same properties.

To prove surjectivity on $\pi_q$, use a finite cube parameter
$P=[0,1]^q$ and put $D=(B\times[0,1])\times P$.
A based class is represented by a definable homeomorphism of $D$ over
$P$ with the bottom fixed and the identity on the parameter boundary.
For $q=0$ the last requirement is empty. Choose a graph formula for
this $\mathcal N$-family. Replace its finitely many parameters outside
$\mathcal M$ by the variable tuple $a$. Impose in $T$ precisely the
total-homeomorphism, parameter, bottom, top and based-boundary
conditions of this family. The original tuple $a$ belongs to one cell
$C(\mathcal N)$ from the partition above.

A definable cell over $\mathcal M$ is definably homeomorphic over
$\mathcal M$ to an open cube of its dimension. Its extension is still
such a cell. A line segment in that open cube gives a definable path
in $C(\mathcal N)$ from the original tuple to the chosen
$a_0\in C(\mathcal M)$. Pull back the jointly continuous evaluation
and inverse evaluation along this path. The result is a homotopy
through exactly the prescribed families; all required boundary
conditions remain true because every tuple lies in $T$.
Its last family is defined over $\mathcal M$.
This proves surjectivity on every based homotopy group and on
components. We use finite parameter families to represent the
homotopy groups of the simplicial realization, as in A.115--A.116.

For injectivity, suppose an $\mathcal M$-family has a based
$\mathcal N$-nullhomotopy. Choose one finite formula for the graph
of that nullhomotopy and make its finitely many extra parameters
variables. The requirement that it be a joint continuous
parameter-preserving homeomorphism with joint continuous inverse,
with the specified fixed subsets, prescribed endpoint graph, and
identity on the other relative faces, is one first-order condition
over $\mathcal M$. It is satisfiable in $\mathcal N$, so
elementarity supplies a solution in $\mathcal M$. This is a
nullhomotopy with the exact same boundary data. The identical
argument descends a path between any two $\mathcal M$ zero-simplices.
Thus extension is injective on all homotopy groups and on components.

Simplicial groups are Kan. The preceding bijections and isomorphisms
therefore prove the asserted weak equivalence. For a finite rational
PL model, extension of a PL map is the same map as extension of its
definable graph. The square commutes literally, and its PL vertical
map is a weak equivalence by A.115. The final assertions follow by
the resulting isomorphisms on every homotopy group.
\end{proof}

\paragraph{Precise corollaries and limitations.}
Fix a finite abstract PL manifold, realized by rational vertices.
For any complete o-minimal theory $T$ with a real model, its whole
PL-to-definable concordance comparison over every model of $T$ has
the same homotopy-degree properties as over that real model. Indeed,
two models of a complete theory embed elementarily into a common
model: the union of their elementary diagrams in disjoint added
constants is consistent by completeness and compactness. The
rational PL model is preserved by both embeddings, so the theorem
and A.115 apply to the resulting comparison zigzag. In particular A.114's
injectivity in every homotopy degree and fundamental-group
isomorphism, when every base component has dimension at least five,
extend to all such models. Surjectivity in degrees at least two is
not proved by this reduction. The real-model
hypothesis cannot simply be discarded: an o-minimal language may,
for example, name a positive infinite constant by its complete
theory, so that no real model exists.

In a countable language one can take a countable elementary submodel
containing the finitely many parameters defining $B$ and the marked
data. There are then only countably many defining formulas and
parameter tuples, giving a countable simplicial model for the same
weak type. This is a cardinality reduction. An arbitrary elementary
diagram is not a decidable theory; no effective enumeration of valid
homeomorphism codes, no oracle-free recursive presentation and no
PL-to-definable weak equivalence follows from it. We do not count
this cardinality observation as a computation closing the original
question.
